%% kapa: exam=1 round=19 year=2008 period=2 form=A subject=economics subject_ko=경제원론 seq=9 number=9
%% source: Q-Net 감정평가사 기출문제 (cst00309) artlSeq=1001962 file=extracted/1008/1001962/2008년 제19회 감정평가사 1차시험 기출문제/2교시.hwp
%% conversion: kordoc 4.24.51 -> tools/gen1.py (auto); review=proofread-sonnet; flags=none
\begin{question}{9}
도덕적 해이에 관한 다음 설명 중 옳은 것을 모두 고르면?
\begin{framedtext}
ㄱ. 본인-대리인(principal-agent) 문제는 도덕적 해이의 한 예이다. \par
ㄴ. 보험 계약에서 사고 피해액의 일정 비율을 가입자가 부담하면 도덕적 해이의 문제가 완화된다. \par
ㄷ. 보험 계약에서 통계적으로 사고가 날 확률이 높은 집단에 속한 사람에게 비싼 보험료를 요구하는 것은 도덕적 해이의 문제를 완화시킨다. \par
ㄹ. 개인의 실적에 따라 성과급을 지급하는 것은 도덕적 해이의 문제를 완화시킨다. \par
\end{framedtext}
\begin{choices}
  \item ㄱ, ㄴ
  \item ㄱ, ㄷ
  \item ㄷ, ㄹ
  \item ㄱ, ㄴ, ㄹ
  \item ㄴ, ㄷ, ㄹ
\end{choices}
\end{question}
\answer{4}
